% !TeX root = ../main.tex

\chapter{Design}

\section{Design Goals}

The goal of this thesis is to design a price manipulation analysis workflow
whose output can be evaluated as an audit report, not only as a binary
vulnerability label. This leads to four design goals.

\textbf{Evidence first.} The system should not ask an LLM to find
vulnerabilities from an entire repository without structure. Static analysis
must first provide concrete source, sink, and path evidence.

\textbf{Semantic grouping.} Raw taint paths are often redundant. A single root
cause may produce many syntactic paths because of SSA variables, helper
functions, tuple returns, or multiple sink operands. The system should group
paths into representative semantic units before LLM analysis.

\textbf{Structured reasoning.} LLM output should be constrained to explicit
fields such as source classification, sink impact decision, root-cause
location, missing mitigation, and supporting evidence IDs. Free-form prose
alone is difficult to validate and score.

\textbf{Audit-oriented evaluation.} The final output should be comparable to
audit findings. The evaluation should measure localization, explanation match,
and over-reporting, not only protocol-level detection.

\section{Workflow Overview}

The proposed workflow contains four major phases:

\begin{enumerate}
  \item Static taint analysis identifies candidate flows from price-related
        sources to economically meaningful sinks.
  \item Path abstraction canonicalizes sources, records provenance, groups
        sink operands, and emits representative paths.
  \item LLM-based analysis filters retained evidence and generates structured
        audit-style findings.
  \item Report-level evaluation matches predictions against audit findings and
        computes binary, localization, semantic, and over-reporting metrics.
\end{enumerate}

This design follows the high-level direction of PMDetector, but changes the
evaluation target. PMDetector-style taint analysis is used as the front end, and
LLM reasoning is used as the semantic filter. The difference is that this
thesis treats the generated report as the main artifact.

\section{Static Analysis Design}

\subsection{Source Model}

The source model intentionally over-approximates possible price-related values.
Candidate sources include oracle-like functions, AMM reserve and price reads,
external view or pure calls with price-like usage, flash-loan-related calls,
and storage reads from price-like state variables. This broad definition is
necessary because DeFi protocols often wrap price reads behind protocol-specific
helper functions.

Each source is represented not only as a raw variable but also as a canonical
source record. A canonical source groups multiple exact seed identifiers that
refer to the same semantic source. For example, a helper function may return a
price that is unpacked into several local variables across different callers.
Treating these as unrelated sources would inflate the number of findings.
Canonicalization allows the system to report one root cause with multiple
supporting paths.

\subsection{Sink Model}

The sink model focuses on economically meaningful operations. Candidate sinks
include price calculations, collateral checks, lending and borrowing operations,
margin updates, liquidation logic, token minting and burning, transfers, and
state writes that affect accounting. The sink model is broader than the final
vulnerability definition because static analysis should first collect possible
evidence.

For evaluation, it is important to record the exact sink operand. A function may
contain several state writes, but only one operand may be controlled by a
manipulable price source. The design therefore tracks semantic sink operands
instead of reporting only that a function contains a sink.

\subsection{Propagation and Path Reconstruction}

Taint propagation operates over Slither IR. The design propagates taint through
assignments, arithmetic operations, references, member accesses, tuple
unpacking, return values, internal calls, library calls, and selected
inter-contract calls. The analysis also records provenance so that a local
wrapper can be connected to the terminal external read.

After propagation, the system reconstructs paths from sink operands back to
their sources. The raw path list is expected to contain duplicates and near
duplicates. This is acceptable at the static-analysis stage because the goal is
to avoid missing potentially relevant flows. Redundancy is reduced in the path
abstraction phase.

\section{Representative Path Design}

The representative path layer converts raw static-analysis results into compact
LLM evidence. It has three responsibilities.

First, it collapses duplicate raw paths that share the same semantic
source-sink relationship. This reduces prompt size and prevents the LLM from
mistaking path volume for vulnerability severity.

Second, it preserves information needed for audit reporting: source label,
source function, source provenance, terminal external reads, sink function,
marked sink lines, state writes, external calls, and path counts. The LLM should
see enough context to reason about exploitability without receiving the entire
repository.

Third, it assigns stable evidence identifiers. The final LLM report refers to
supporting evidence IDs, making it possible to trace each generated finding
back to concrete static-analysis artifacts.

\section{LLM Analysis Design}

\subsection{Stage 0: Source Classification}

The first LLM stage classifies each canonical source. The allowed
classifications are \texttt{spot\_amm\_price},
\texttt{twap\_or\_oracle\_price}, \texttt{protocol\_state},
\texttt{fee\_or\_liquidity}, \texttt{unrelated}, and \texttt{unknown}. The
model also outputs a decision: keep, drop, or need more context.

This stage is designed to remove obvious false positives early. For example,
pool addresses, DAO addresses, balances, total supply, and fee-growth helpers
may be part of a taint path but are not necessarily manipulable price sources.

\subsection{Stage 1: Sink Group Triage}

For retained sources, the second LLM stage judges whether each sink group is
economically dangerous. A sink is dangerous when the tainted value affects
accounting, margin, collateral, liquidation, value transfer, minting, burning,
or another price-sensitive operation. It is not dangerous when it is telemetry,
metadata, harmless caching, or a value not actually controlled by the retained
source.

\subsection{Stage 2: Root-Cause Judgment}

The final LLM stage emits at most one finding per canonical source. This design
reduces over-reporting by preventing every path from becoming a separate
finding. Each output includes a verdict, root-cause location, vulnerable
external read, missing mitigation, supporting evidence IDs, duplicate or
supporting evidence IDs, and a concise explanation.

The output schema is intentionally audit-like. It should be possible to compare
the generated explanation with a ground-truth audit finding and ask whether
they describe the same issue.

\section{Evaluation Design}

\subsection{Ground Truth}

The proposed ground truth is a curated set of audit findings related to price
manipulation. Each finding should contain at least the protocol, title,
description, affected function or code location when available, impact, and
mitigation. The dataset should also record whether source code is available and
whether the repository is build-verified.

\subsection{Report Matching}

For each project, generated findings are matched against ground-truth audit
findings. A pair score combines localization agreement, mitigation agreement,
and semantic explanation similarity. Low-scoring pairs are treated as
non-matches. Maximum-weight bipartite matching can then be used to assign
predictions to ground truth without allowing one prediction to satisfy multiple
findings.

\subsection{Over-Reporting Penalty}

The evaluation penalizes unsupported extra findings. For a project $p$, the
basic penalty is:

\[
\mathrm{Penalty}(p) = \max(0, \#\mathrm{Predictions}(p) -
\#\mathrm{GroundTruth}(p)).
\]

This penalty captures audit workload. A detector that emits many weak reports
can increase recall, but it also increases manual review cost.

\section{Scope}

The initial scope is price manipulation in Solidity-based DeFi protocols. The
dataset may contain related reward or vote manipulation findings, but the core
thesis focuses on price-related sources and price-sensitive sinks. Other
vulnerability classes such as reentrancy, access control, and bridge-specific
logic are outside the main scope unless they directly affect price manipulation
mitigation.
